% --- Archon physics formalization source begin ---
% archon:physics
% archon:covers PhyXMiniProblems/problem_phyx_mini_0322.lean
% archon:source-report reports/phyx_mini/problem_phyx_mini_0322.source.json

\chapter{Physics problem phyx\_mini\_0322}
\label{ch:PhyXMiniProblems_problem_phyx_mini_0322}

\paragraph{Problem source.}
\#\# Physical scenario

In figure, sound waves \$A\$ and \$B\$, both of wavelength \$\textbackslash{}lambda\$, are initially in phase and traveling rightward, as indicated by the two rays. Wave \$A\$ is reflected from four surfaces but ends up traveling in its original direction.

\#\# Question

What multiple of wavelength \$\textbackslash{}lambda\$ is the smallest value of distance \$L\$ in the figure that puts \$A\$ and \$B\$ exactly out of phase with each other after the reflections?

\#\# Answer choices

- A: A: \$\textbackslash{}lambda\$

- B: B: \$0.75 \textbackslash{}lambda\$

- C: C: \$0.5 \textbackslash{}lambda\$

- D: D: \$0.25 \textbackslash{}lambda\$

\#\# Figure caption (auxiliary; use the image as primary evidence)

The image shows a diagram with a series of red arrows connected in a zigzag pattern forming a descending staircase-like structure on the left. These arrows are labeled "A" at the top left and are connected by horizontal and vertical segments. The vertical segments have a consistent length labeled as "L." Blue rectangles are positioned diagonally at each turn in the zigzag pattern.

To the lower part of the image, there is a straight horizontal red arrow labeled "B," running parallel to the base of the zigzag structure. Both "A" and "B" arrows point from left to right.

The blue rectangles and red arrows suggest a directional or flow process from point A to B, with the zigzag path representing a more complex or resistant route compared to the direct path represented by B.

\#\# Dataset metadata

Wave Properties; Spatial Relation Reasoning, Implicit Condition Reasoning

\paragraph{Recorded answer/context.}
D: D: \$0.25 \textbackslash{}lambda\$

\paragraph{Figure/image path.}
/root/proposal\_for\_physic/science-mango/phyx\_mini\_run/phyx\_data/test\_image/322.png

\paragraph{Formalization target.}
create a compiling Lean file with sorry bodies at `PhyXMiniProblems/problem\_phyx\_mini\_0322.lean`. The Lean declarations must preserve the physical quantities, dimensions or dimensional roles, figure labels, governing-law hypotheses, and final relation expressed by this problem.
Use Mathlib/Physlib names found through LeanExplore where available. If a physics API is missing, introduce faithful local abstractions rather than scalar placeholder aliases.

\begin{theorem}[Physics formalization target]
\label{thm:physics:phyx_mini_0322:target}
\lean{PhyXMiniProblems.ProblemPhyXMini0322.problem_phyx_mini_0322}
\uses{def:physics:phyx-mini-0322:phyxminiproblems-problemphyxmini0322-lengthreadout, def:physics:phyx-mini-0322:phyxminiproblems-problemphyxmini0322-fourreflectionsoundsetup, def:physics:phyx-mini-0322:phyxminiproblems-problemphyxmini0322-matchesproblemandfigure, def:physics:phyx-mini-0322:phyxminiproblems-problemphyxmini0322-hasphysicalacousticparameters, def:physics:phyx-mini-0322:phyxminiproblems-problemphyxmini0322-satisfiesdepictedpathgeometry, def:physics:phyx-mini-0322:phyxminiproblems-problemphyxmini0322-satisfiesuniformreflectionphaselaw, def:physics:phyx-mini-0322:phyxminiproblems-problemphyxmini0322-satisfiesacousticphaseaccumulationlaw, def:physics:phyx-mini-0322:phyxminiproblems-problemphyxmini0322-isleastpositiveoutofphaseleglength, def:physics:phyx-mini-0322:phyxminiproblems-problemphyxmini0322-displayedwavelengthmultiple, def:physics:phyx-mini-0322:phyxminiproblems-problemphyxmini0322-isuniquematchinganswer, def:physics:phyx-mini-0322:phyxminiproblems-problemphyxmini0322-recordeddatasetanswer}
The assigned autoformalize agent should translate this physics problem into Lean declarations in the covered file, with theorem and lemma proof bodies written as `by sorry`.
\end{theorem}
\begin{proof}
This is an autoformalization task, not a proof task. Produce statements that can later be proved by the physics prover without weakening the physical model.
\end{proof}

% --- Archon declaration topology begin ---
\section{Declaration topology}
\begin{definition}[Acoustic Length]
  \label{def:physics:phyx-mini-0322:phyxminiproblems-problemphyxmini0322-acousticlength}
  \lean{PhyXMiniProblems.ProblemPhyXMini0322.AcousticLength}
  A signed physical length, independent of the unit used to read it
\end{definition}
\begin{proof}
  This is the declared model definition of Acoustic Length.
\end{proof}
\begin{definition}[length Readout]
  \label{def:physics:phyx-mini-0322:phyxminiproblems-problemphyxmini0322-lengthreadout}
  \lean{PhyXMiniProblems.ProblemPhyXMini0322.lengthReadout}
  \uses{def:physics:phyx-mini-0322:phyxminiproblems-problemphyxmini0322-acousticlength}
  Read a physical length as a real scalar in a selected length unit
\end{definition}
\begin{proof}
  This is the declared model definition of length Readout.
\end{proof}
\begin{definition}[length In Meters]
  \label{def:physics:phyx-mini-0322:phyxminiproblems-problemphyxmini0322-lengthinmeters}
  \lean{PhyXMiniProblems.ProblemPhyXMini0322.lengthInMeters}
  \uses{def:physics:phyx-mini-0322:phyxminiproblems-problemphyxmini0322-acousticlength, def:physics:phyx-mini-0322:phyxminiproblems-problemphyxmini0322-lengthreadout}
  Metre readout of a physical length
\end{definition}
\begin{proof}
  This is the declared model definition of length In Meters.
\end{proof}
\begin{definition}[length Of Meters]
  \label{def:physics:phyx-mini-0322:phyxminiproblems-problemphyxmini0322-lengthofmeters}
  \lean{PhyXMiniProblems.ProblemPhyXMini0322.lengthOfMeters}
  \uses{def:physics:phyx-mini-0322:phyxminiproblems-problemphyxmini0322-acousticlength}
  Construct a physical length from its metre readout
\end{definition}
\begin{proof}
  This is the declared model definition of length Of Meters.
\end{proof}
\begin{definition}[Wave Label]
  \label{def:physics:phyx-mini-0322:phyxminiproblems-problemphyxmini0322-wavelabel}
  \lean{PhyXMiniProblems.ProblemPhyXMini0322.WaveLabel}
  The two sound-wave rays labelled in the primary image
\end{definition}
\begin{proof}
  This is the declared model definition of Wave Label.
\end{proof}
\begin{definition}[Horizontal Direction]
  \label{def:physics:phyx-mini-0322:phyxminiproblems-problemphyxmini0322-horizontaldirection}
  \lean{PhyXMiniProblems.ProblemPhyXMini0322.HorizontalDirection}
  Horizontal propagation directions used by the problem statement
\end{definition}
\begin{proof}
  This is the declared model definition of Horizontal Direction.
\end{proof}
\begin{definition}[Reflection Surface]
  \label{def:physics:phyx-mini-0322:phyxminiproblems-problemphyxmini0322-reflectionsurface}
  \lean{PhyXMiniProblems.ProblemPhyXMini0322.ReflectionSurface}
  The four blue reflecting surfaces, named in the order in which wave A encounters them while descending through the staircase
\end{definition}
\begin{proof}
  This is the declared model definition of Reflection Surface.
\end{proof}
\begin{definition}[Four Reflection Sound Setup]
  \label{def:physics:phyx-mini-0322:phyxminiproblems-problemphyxmini0322-fourreflectionsoundsetup}
  \lean{PhyXMiniProblems.ProblemPhyXMini0322.FourReflectionSoundSetup}
  \uses{def:physics:phyx-mini-0322:phyxminiproblems-problemphyxmini0322-acousticlength, def:physics:phyx-mini-0322:phyxminiproblems-problemphyxmini0322-wavelabel, def:physics:phyx-mini-0322:phyxminiproblems-problemphyxmini0322-horizontaldirection, def:physics:phyx-mini-0322:phyxminiproblems-problemphyxmini0322-reflectionsurface}
  The physical quantities, ray data, and phase observables of the diagram. totalPathLengthAtLegLength wave candidateL describes the member of the depicted family obtained by assigning candidateL to both vertical legs. Thus the observable wavesExactlyOutOfPhaseAtLegLength is independent data constrained by the general acoustic phase laws below; it is not defined to select the recorded answer
\end{definition}
\begin{proof}
  This is the declared model definition of Four Reflection Sound Setup.
\end{proof}
\begin{definition}[Matches Problem And Figure]
  \label{def:physics:phyx-mini-0322:phyxminiproblems-problemphyxmini0322-matchesproblemandfigure}
  \lean{PhyXMiniProblems.ProblemPhyXMini0322.MatchesProblemAndFigure}
  \uses{def:physics:phyx-mini-0322:phyxminiproblems-problemphyxmini0322-fourreflectionsoundsetup}
  Qualitative and integer data stated in the problem and read from the image. The two waves begin in phase and rightward, wave A meets the four displayed surfaces in order, wave B meets none, and both outgoing rays point right. The raster shows exactly two vertical legs labelled L. No wavelength multiple or answer choice occurs in these data
\end{definition}
\begin{proof}
  This is the declared model definition of Matches Problem And Figure.
\end{proof}
\begin{definition}[Has Physical Acoustic Parameters]
  \label{def:physics:phyx-mini-0322:phyxminiproblems-problemphyxmini0322-hasphysicalacousticparameters}
  \lean{PhyXMiniProblems.ProblemPhyXMini0322.HasPhysicalAcousticParameters}
  \uses{def:physics:phyx-mini-0322:phyxminiproblems-problemphyxmini0322-lengthreadout, def:physics:phyx-mini-0322:phyxminiproblems-problemphyxmini0322-fourreflectionsoundsetup}
  Positivity of the common wavelength and of the representative L drawn in the setup. Positivity for alternative candidate leg lengths is imposed in the least-positive predicate below
\end{definition}
\begin{proof}
  This is the declared model definition of Has Physical Acoustic Parameters.
\end{proof}
\begin{definition}[Satisfies Depicted Path Geometry]
  \label{def:physics:phyx-mini-0322:phyxminiproblems-problemphyxmini0322-satisfiesdepictedpathgeometry}
  \lean{PhyXMiniProblems.ProblemPhyXMini0322.SatisfiesDepictedPathGeometry}
  \uses{def:physics:phyx-mini-0322:phyxminiproblems-problemphyxmini0322-acousticlength, def:physics:phyx-mini-0322:phyxminiproblems-problemphyxmini0322-lengthreadout, def:physics:phyx-mini-0322:phyxminiproblems-problemphyxmini0322-fourreflectionsoundsetup}
  The horizontal pieces of the two rays have the same total horizontal extent. Consequently wave A exceeds the straight path of wave B precisely by its two vertical legs. The relation is stated for every proposed L and every length unit, and does not single out a numerical multiple of lambda
\end{definition}
\begin{proof}
  This is the declared model definition of Satisfies Depicted Path Geometry.
\end{proof}
\begin{definition}[Satisfies Uniform Reflection Phase Law]
  \label{def:physics:phyx-mini-0322:phyxminiproblems-problemphyxmini0322-satisfiesuniformreflectionphaselaw}
  \lean{PhyXMiniProblems.ProblemPhyXMini0322.SatisfiesUniformReflectionPhaseLaw}
  \uses{def:physics:phyx-mini-0322:phyxminiproblems-problemphyxmini0322-wavelabel, def:physics:phyx-mini-0322:phyxminiproblems-problemphyxmini0322-reflectionsurface, def:physics:phyx-mini-0322:phyxminiproblems-problemphyxmini0322-fourreflectionsoundsetup}
  A reflection from any one of the four like surfaces has the same conventional phase response. Depending on whether acoustic pressure or displacement is tracked, that common response is either zero or one half-cycle. The total reflection contribution is the sum along the ray's displayed reflection sequence. In either convention, four identical reflections contribute a whole number of cycles; this law contains no condition on L
\end{definition}
\begin{proof}
  This is the declared model definition of Satisfies Uniform Reflection Phase Law.
\end{proof}
\begin{definition}[Satisfies Acoustic Phase Accumulation Law]
  \label{def:physics:phyx-mini-0322:phyxminiproblems-problemphyxmini0322-satisfiesacousticphaseaccumulationlaw}
  \lean{PhyXMiniProblems.ProblemPhyXMini0322.SatisfiesAcousticPhaseAccumulationLaw}
  \uses{def:physics:phyx-mini-0322:phyxminiproblems-problemphyxmini0322-acousticlength, def:physics:phyx-mini-0322:phyxminiproblems-problemphyxmini0322-lengthreadout, def:physics:phyx-mini-0322:phyxminiproblems-problemphyxmini0322-fourreflectionsoundsetup}
  Propagation over one extra wavelength delays a wave by one phase cycle. Therefore the arrival phase of A relative to B is its initial relative phase, minus its excess path divided by lambda, plus the difference of the two accumulated reflection phases. Exact opposition means an odd half-cycle modulo an integral number of cycles. These are general phase laws quantified over arbitrary proposed leg lengths; they do not assert which proposed length is the least solution
\end{definition}
\begin{proof}
  This is the declared model definition of Satisfies Acoustic Phase Accumulation Law.
\end{proof}
\begin{definition}[Is Least Positive Out Of Phase Leg Length]
  \label{def:physics:phyx-mini-0322:phyxminiproblems-problemphyxmini0322-isleastpositiveoutofphaseleglength}
  \lean{PhyXMiniProblems.ProblemPhyXMini0322.IsLeastPositiveOutOfPhaseLegLength}
  \uses{def:physics:phyx-mini-0322:phyxminiproblems-problemphyxmini0322-acousticlength, def:physics:phyx-mini-0322:phyxminiproblems-problemphyxmini0322-lengthinmeters, def:physics:phyx-mini-0322:phyxminiproblems-problemphyxmini0322-fourreflectionsoundsetup}
  A physical leg length is the requested one when it is positive, puts the two waves exactly out of phase, and has no larger metre readout than any other positive leg length doing so. Metre readout is used only to express the order on physical lengths
\end{definition}
\begin{proof}
  This is the declared model definition of Is Least Positive Out Of Phase Leg Length.
\end{proof}
\begin{definition}[Answer Choice]
  \label{def:physics:phyx-mini-0322:phyxminiproblems-problemphyxmini0322-answerchoice}
  \lean{PhyXMiniProblems.ProblemPhyXMini0322.AnswerChoice}
  Labels of the four printed answer choices
\end{definition}
\begin{proof}
  This is the declared model definition of Answer Choice.
\end{proof}
\begin{definition}[displayed Wavelength Multiple]
  \label{def:physics:phyx-mini-0322:phyxminiproblems-problemphyxmini0322-displayedwavelengthmultiple}
  \lean{PhyXMiniProblems.ProblemPhyXMini0322.displayedWavelengthMultiple}
  \uses{def:physics:phyx-mini-0322:phyxminiproblems-problemphyxmini0322-answerchoice}
  The multiple of the common wavelength printed beside each answer label
\end{definition}
\begin{proof}
  This is the declared model definition of displayed Wavelength Multiple.
\end{proof}
\begin{definition}[Matches Displayed Wavelength Multiple]
  \label{def:physics:phyx-mini-0322:phyxminiproblems-problemphyxmini0322-matchesdisplayedwavelengthmultiple}
  \lean{PhyXMiniProblems.ProblemPhyXMini0322.MatchesDisplayedWavelengthMultiple}
  \uses{def:physics:phyx-mini-0322:phyxminiproblems-problemphyxmini0322-acousticlength, def:physics:phyx-mini-0322:phyxminiproblems-problemphyxmini0322-lengthreadout, def:physics:phyx-mini-0322:phyxminiproblems-problemphyxmini0322-fourreflectionsoundsetup, def:physics:phyx-mini-0322:phyxminiproblems-problemphyxmini0322-isleastpositiveoutofphaseleglength, def:physics:phyx-mini-0322:phyxminiproblems-problemphyxmini0322-answerchoice, def:physics:phyx-mini-0322:phyxminiproblems-problemphyxmini0322-displayedwavelengthmultiple}
  A displayed multiple matches the physical model when some least positive out-of-phase leg length has that ratio to the wavelength in every unit
\end{definition}
\begin{proof}
  This is the declared model definition of Matches Displayed Wavelength Multiple.
\end{proof}
\begin{definition}[Is Unique Matching Answer]
  \label{def:physics:phyx-mini-0322:phyxminiproblems-problemphyxmini0322-isuniquematchinganswer}
  \lean{PhyXMiniProblems.ProblemPhyXMini0322.IsUniqueMatchingAnswer}
  \uses{def:physics:phyx-mini-0322:phyxminiproblems-problemphyxmini0322-fourreflectionsoundsetup, def:physics:phyx-mini-0322:phyxminiproblems-problemphyxmini0322-answerchoice, def:physics:phyx-mini-0322:phyxminiproblems-problemphyxmini0322-matchesdisplayedwavelengthmultiple}
  Exactly one printed choice agrees with the least positive physical length
\end{definition}
\begin{proof}
  This is the declared model definition of Is Unique Matching Answer.
\end{proof}
\begin{definition}[recorded Dataset Answer]
  \label{def:physics:phyx-mini-0322:phyxminiproblems-problemphyxmini0322-recordeddatasetanswer}
  \lean{PhyXMiniProblems.ProblemPhyXMini0322.recordedDatasetAnswer}
  \uses{def:physics:phyx-mini-0322:phyxminiproblems-problemphyxmini0322-answerchoice}
  The answer label recorded in the supplied dataset metadata
\end{definition}
\begin{proof}
  This is the declared model definition of recorded Dataset Answer.
\end{proof}
\begin{lemma}[least Positive Out Of Phase Leg Length is quarter Wavelength]
  \label{lem:physics:phyx-mini-0322:phyxminiproblems-problemphyxmini0322-leastpositiveoutofphaseleglength-is-quarterwavelength}
  \lean{PhyXMiniProblems.ProblemPhyXMini0322.leastPositiveOutOfPhaseLegLength_is_quarterWavelength}
  \uses{def:physics:phyx-mini-0322:phyxminiproblems-problemphyxmini0322-lengthreadout, def:physics:phyx-mini-0322:phyxminiproblems-problemphyxmini0322-fourreflectionsoundsetup, def:physics:phyx-mini-0322:phyxminiproblems-problemphyxmini0322-matchesproblemandfigure, def:physics:phyx-mini-0322:phyxminiproblems-problemphyxmini0322-hasphysicalacousticparameters, def:physics:phyx-mini-0322:phyxminiproblems-problemphyxmini0322-satisfiesdepictedpathgeometry, def:physics:phyx-mini-0322:phyxminiproblems-problemphyxmini0322-satisfiesuniformreflectionphaselaw, def:physics:phyx-mini-0322:phyxminiproblems-problemphyxmini0322-satisfiesacousticphaseaccumulationlaw, def:physics:phyx-mini-0322:phyxminiproblems-problemphyxmini0322-isleastpositiveoutofphaseleglength}
  The least positive out-of-phase leg length is one quarter of the common wavelength. This relation is unit independent and identifies the displayed answer D. It is a derived conclusion, not a field of the setup or of any law
\end{lemma}
\begin{proof}
  The conclusion follows from the governing relations and figure readouts named in the statement of least Positive Out Of Phase Leg Length is quarter Wavelength.
\end{proof}
% --- Archon declaration topology end ---
% --- Archon physics formalization source end ---
